\section{Operator-aware terminal contraction for finite mixtures}
\label{sec:terminal_history}
This section retains the original RED update and permits any finite
number of equal-width isotropic Gaussian components, any ambient
dimension, and any linear observation operator.
Write $m_v=\sum_iw_i\mu_i$ and
$C_v=\sum_iw_i(\mu_i-m_v)(\mu_i-m_v)^\top$.
This is the posterior covariance of the template means.
It excludes the within-component covariance.
Throughout, $v>\tau^2$, $c=\lambda(1-\tau^2/v)>0$ and $B=A^\top A$.
Conditional mean--covariance derivative identities have an established
foundation~\cite{dytso2023}. Our use is a model-input certificate for
the original RED terminal map.

\subsection{An exact local test using the sensing operator}
Let $\mathcal S=\operatorname{span}\{\mu_i-\mu_0\}$ and let
$U,V$ be orthonormal bases of $\mathcal S$ and its orthogonal
complement. Define
\begin{align}
 G_c={}&U^\top BU+cI \nonumber\\
 &-U^\top BV(V^\top BV+cI)^{-1}V^\top BU,\\
 C_S&=U^\top C_vU,\qquad
 \rho_v(x)=\frac cv\lambda_{\max}(G_c^{-1/2}C_SG_c^{-1/2}).
\end{align}
The term involving $V$ vanishes when $\mathcal S$ is the whole space.
For coincident means set $\rho_v=0$.
\begin{proposition}[Operator--posterior curvature test]\label{prop:operator_posterior}
The Hessian of the exact RED potential is
\[
 H_v=B+cI-(c/v)C_v.
\]
At each state, $H_v\succ0$ if and only if $\rho_v(x)<1$.
On a convex set where $\rho_v\le\bar\rho<1$,
\[
 (1-\bar\rho)(\lambda_{\min}(B)+c)I
 \preceq H_v\preceq(\lambda_{\max}(B)+c)I.
\]
For $0<\gamma\le(\lambda_{\max}(B)+c)^{-1}$,
the map is Lipschitz with factor
$q=1-\gamma(1-\bar\rho)(\lambda_{\min}(B)+c)<1$ on that set.
\end{proposition}
\begin{proof}
Differentiating the responsibilities gives $D_xm_v=C_v/v$.
In the basis $[U,V]$, the lower-right block of $H_v$ is positive
definite and its Schur complement is $G_c-(c/v)C_S$.
Moreover $U^\top(B+cI)^{-1}U=G_c^{-1}$, so the nonzero eigenvalues
in the definition of $\rho_v$ equal those of
$(c/v)(B+cI)^{-1/2}C_v(B+cI)^{-1/2}$.
The Loewner bounds follow.
Integrating $I-\gamma H_v$ on a line segment in the convex set
gives the stated Lipschitz factor.
\end{proof}
Constructing $G_c$ requires solves involving the actual operator.
Only the final covariance eigenproblem has dimension at most $K-1$.
Pairwise secant gains alone do not lower-bound $B$ on $\mathcal S$.
For example, means $(0,0),(1,0),(2,1)$ under
$B=\operatorname{diag}(1,0)$ have pair gains $1,4/5,1/2$ but
span a null direction.

\subsection{A common invariant terminal ball}
Fix the terminal map and a ball
$\mathcal B=\{x:\norm{x-\bar x}\le R\}$.
For reference component $j$, set $d_{ij}=\mu_i-\mu_j$ and
\begin{align}
 \Delta_{ij}&=\norm{\bar x-\mu_i}^2-\norm{\bar x-\mu_j}^2,\\
 \beta_{ij}&=\operatorname{logistic}\left(
 \log\frac{\pi_i}{\pi_j}-\frac{\Delta_{ij}-2R\norm{d_{ij}}}{2v_L}\right),\\
 \overline C_{\mathcal B}&=\sum_{i\ne j}\beta_{ij}d_{ij}d_{ij}^\top.
\end{align}
The affine distance-gap bound implies $w_i/w_j\le
\exp(\operatorname{logit}\beta_{ij})$, hence $w_i\le\beta_{ij}$.
Since
$C_v=\sum_iw_id_{ij}d_{ij}^\top-(m_v-\mu_j)(m_v-\mu_j)^\top$,
this is a valid covariance upper bound on the whole ball.

\begin{proposition}[Common terminal ball]\label{prop:common_terminal_ball}
Let $\overline C$ be any whole-ball covariance upper bound and put
$M=B+c_LI$. Suppose $(c_L/v_L)\overline C\preceq\bar\rho M$
for some $\bar\rho<1$, the step condition of
Proposition~\ref{prop:operator_posterior} holds, and
\[
 \norm{T_{v_L}(\bar x)-\bar x}\le(1-q)R.
\]
Then the ball is invariant and contains a unique terminal fixed point.
For any two starts in the ball at common time $H$,
\[
 \norm{x^{(a)}_{H+n}-x^{(b)}_{H+n}}
 \le q^n\norm{x^{(a)}_H-x^{(b)}_H},\qquad n\ge0.
\]
\end{proposition}
\begin{proof}
Contraction gives
$\norm{T(x)-\bar x}\le q\norm{x-\bar x}+\norm{T(\bar x)-\bar x}\le R$.
The contraction mapping theorem applies on the closed ball.
Iteration of the Lipschitz estimate proves the distance bound.
\end{proof}
If both updates have norm errors bounded by $\delta_a,\delta_b$
and remain in the ball, the distance bound gains
$(\delta_a+\delta_b)(1-q^n)/(1-q)$.
For perturbed invariance add the corresponding error bound to the
center residual.
An almost-sure certificate with a uniform $q$ implies
\[
 |\sqrt{R_a(n)}-\sqrt{R_b(n)}|
 \le q^n\big(\mathbb E\norm{x_H^{(a)}-x_H^{(b)}}^2\big)^{1/2}.
\]
A good-event-only certificate also needs an exceptional-event
$L^2$ term. A finite batch does not imply a population risk guarantee.
Neither common convergence nor a unique ball fixed point implies
accurate recovery.

\subsection{Sharper pairwise and global covariance bounds}
For each pair let
\begin{align}
 t_{ij}&=\log(\pi_i/\pi_j)
 -\frac{\norm{\bar x-\mu_i}^2-\norm{\bar x-\mu_j}^2}{2v},\\
 L_{ij}&=t_{ij}-R\norm{\mu_i-\mu_j}/v,\nonumber\\
 U_{ij}&=t_{ij}+R\norm{\mu_i-\mu_j}/v .
\end{align}
Let $\alpha_i=(\sum_k e^{L_{ki}})^{-1}$ and let $t^\circ_{ij}$
be the projection of zero onto $[L_{ij},U_{ij}]$.
Every state in the ball has $w_i\le\alpha_i$ and
\[
 w_iw_j\le\kappa_{ij}
 =\min\left\{\alpha_i\alpha_j,
 \frac1{(1+e^{t^\circ_{ij}})(1+e^{-t^\circ_{ij}})}\right\}.
\]
The second bound follows by maximizing the product at fixed pair
odds and total mass at most one.
The exact identity
$C_v=\sum_{i<j}w_iw_j(\mu_i-\mu_j)(\mu_i-\mu_j)^\top$
then gives a whole-ball bound by replacing $w_iw_j$ with $\kappa_{ij}$.
One may choose the valid covariance envelope with smaller normalized
spectral bound. An elementwise minimum of matrices is unjustified.

A separate global sufficient bound uses the whitened diameter
$D_M^2=\max_{i,j}(\mu_i-\mu_j)^\top M^{-1}(\mu_i-\mu_j)$.
The directional bounded-variable variance inequality yields
$\rho_v(x)\le cD_M^2/(4v)$ everywhere.
If this bound is below one, the terminal map is globally contractive
at the prescribed step size.
Failing a sufficient bound leaves the global question open.

\subsection{Certificate for the fixed patch batch}
The template span in the photo-patch case has dimension seven.
Its restricted sensing spectrum has minimum $.02509617$,
compared with the pairwise secant minimum $.22752$.
The whitened-diameter global bound is $1.21022$, so it is inconclusive.
For each observation we compute a candidate center by ten Newton
steps from its 5000-call terminal state.
At common time 154 the radius is 1.01 times the maximum distance
of the three saved outputs from this center, with floor $10^{-9}$.
No truth patch is used in this construction.
The original reference-component envelope passes 255 of 256 balls.
The pairwise envelope passes all 256.

The program \path{experiments/code/patch_terminal_interval.py}
uses 60-digit outward-rounded interval arithmetic to check all 256.
It verifies the Gram decomposition with exact rational arithmetic,
proves positive definiteness of
$.45(B+c_LI)-(c_L/v_L)\overline C_{\mathcal B}$ by interval LDL,
and checks center residuals, coverage and step size.
Since $B\succeq0$ and $\norm{B}=1$, it yields the uniform factor
$q=1-\gamma(.55c_L)<.999429$.
The certificate treats exported binary64 coefficients and snapshots
as exact inputs to the subsequent exact terminal map.
It does not enclose the earlier floating-point propagation or
subsequent finite-precision perturbations.
It proves neither a new-observation risk bound nor low error of the
shared fixed points.

\subsection{Related contraction and recovery results}
Sinha et al.~\cite{sinha2025} study fixed linear kernel denoisers,
including appropriately scaled PnP for nonsymmetric kernels.
The present covariance varies with state and the update is residual-additive RED.
Wolf et al.~\cite{wolf2026} jointly analyze denoising covariance
and generalized backprojection in PnP forward--backward splitting.
Their pointwise and distributional bounds use support and restricted
conditioning assumptions. Their compact-support assumptions and
update do not directly apply to our untruncated GMM RED model.
The covariance identities, Schur complements and contraction tools
are established mathematics. The present statements organize them
as model-input checks and a batch-specific terminal certificate.

\section{Coefficient and posterior changes in a scale update}
\label{sec:coefficient_split}
Write $T_{c,v}(x)=x-\gamma[Bx-b+c(x-m_v(x))]$.
The original algorithm remains on $c=\lambda(1-\tau^2/v)$.
For fixed $x$ the exact identity is
\begin{align}
 T_{c,v}(x)-T_{\tilde c,\tilde v}(x)
 ={}&-\gamma(c-\tilde c)(x-m_v(x))\nonumber\\
 &+\gamma\tilde c(m_v(x)-m_{\tilde v}(x)).
\end{align}
For $z=1/v$ and $E_i(x)=\norm{x-\mu_i}^2/2$,
\begin{align*}
 \partial_z m_{1/z}(x)&=-\operatorname{Cov}_{w}(\mu,E),\\
 \norm{\operatorname{Cov}_w(\mu,E)}
 &\le\sqrt{\lambda_{\max}(C_w)\operatorname{Var}_w(E)} .
\end{align*}
Both formulas follow by differentiating the normalized exponential
weights and applying directional Cauchy--Schwarz.
Integration bounds $\norm{m_v-m_{\tilde v}}$ by
$|v^{-1}-\tilde v^{-1}|$ times the supremum of this covariance norm.
On a region containing both trajectories, if the baseline map is
$q_k$-Lipschitz, let
$R_k=\sup_x\norm{x-m_{v_k}(x)}$ and
$M_k=\sup_{x,z}\norm{\operatorname{Cov}_{w(x,1/z)}(\mu,E(x))}$.
Then their distance satisfies
\begin{align}
 A_\ell&=|c_\ell-\tilde c_\ell|R_\ell+
 \tilde c_\ell|v_\ell^{-1}-\tilde v_\ell^{-1}|M_\ell,\nonumber\\
 e_n&\le\left(\prod_{k=0}^{n-1}q_k\right)e_0
 +\gamma\sum_{\ell=0}^{n-1}
 \left(\prod_{k=\ell+1}^{n-1}q_k\right)A_\ell .
\end{align}
This is the iterated one-step comparison.
It permits cancellation of vector effects but does not identify
unique statistical contributions from a single measured output difference.

There is an explicit interaction when two separate ablations are
both measured from the corner $(c_0,v_0)$.
For $m_i=m_{v_i}(x)$,
\begin{align}
 &T_{c_1,v_1}-T_{c_1,v_0}-T_{c_0,v_1}+T_{c_0,v_0}\nonumber\\
 &\hspace{18mm}=\gamma(c_1-c_0)(m_1-m_0).
\end{align}
Thus coefficient-only and posterior-only changes do not generally
sum to the combined change. Its interaction norm is at most
$\gamma|c_1-c_0|\norm{m_1-m_0}$, before state propagation.
